% Part: sets-functions-relations
% Chapter: sets
% Section: important-sets

\documentclass[../../../include/open-logic-section]{subfiles}

\begin{document}

\olfileid[pt]{sfr}{set}{imp}
\olsection{Alguns Conjuntos Importantes}

\begin{ex}
Trabalharemos sobretudo com conjuntos cujos !!{element}s são
objetos matemáticos. Quatro desses conjuntos são tão importantes que têm
nomes específicos:
\begin{multline*}
    \Nat = \{0, 1, 2, 3, \ldots\} \\
    \shoveright{\text{o conjunto dos números naturais}}\\
    \shoveleft{\Int = \{\ldots, -2, -1, 0, 1, 2, \ldots\}} \\
    \shoveright{\text{o conjunto dos números inteiros}}\\
    \shoveleft{\Rat = \Setabs{\nicefrac{m}{n}}{m, n \in \Int\text{ e }n \neq 0}}\\
    \shoveright{\text{o conjunto dos números racionais}}\\
    \shoveleft{\Real = (-\infty, \infty)}\\
    \text{o conjunto dos números reais (o contínuo)}
\end{multline*}
São todos conjuntos \emph{infinitos}, isto é, cada um tem
infinitos !!{element}s.

À medida que percorremos estes conjuntos, vamos acrescentando \emph{mais} números ao
nosso repertório. De fato, deve ficar claro que $\Nat \subseteq \Int
\subseteq \Rat \subseteq \Real$: afinal, todo número natural é um
inteiro; todo inteiro é racional; e todo racional é real.
Igualmente, deve ficar claro que $\Nat \subsetneq \Int \subsetneq
\Rat$, pois $-1$ é inteiro mas não natural, e
$\nicefrac{1}{2}$ é racional mas não inteiro. É menos óbvio
que $\Rat \subsetneq \Real$, isto é, que há números reais
que não são racionais\oliflabeldef{sfr:arith:real:realline}{, mas voltaremos a isso na \olref[arith][real]{realline}}{}. 

Também usaremos por vezes o conjunto dos inteiros positivos $\PosInt = \{1,
2, 3, \dots\}$ e o conjunto que contém apenas os dois primeiros naturais
$\Bin = \{0, 1\}$.
\end{ex}

\begin{tagblock}{compsci}
\begin{ex}[Cadeias] 
Outro exemplo interessante é o conjunto $A^{*}$ das \emph{cadeias finitas}
sobre um alfabeto $A$: qualquer sequência finita de elementos de~$A$
é uma cadeia sobre $A$. Incluímos a \emph{cadeia vazia $\Lambda$}
entre as cadeias sobre~$A$, para todo alfabeto~$A$. Por exemplo,
\begin{multline*}
\Bin^*
=\{\Lambda,0,1,00,01,10,11,\\
000,001,010,011,100,101,110,111,0000,\ldots\}.
\end{multline*}
Se $x=x_{1}\ldots x_{n}\in A^{*}$ é uma cadeia constituída por $n$
``letras'' de $A$, então dizemos que o \emph{comprimento} da cadeia é~$n$
e escrevemos $\len{x}=n$.
\end{ex}
\end{tagblock}

\begin{ex}[Sequências infinitas]
Para qualquer conjunto $A$ podemos também considerar o conjunto~$A^\omega$ das sequências infinitas
de !!{element}s de~$A$. Uma sequência infinita
$a_1a_2a_3a_4\dots$ consiste numa lista infinita numa só direção de objetos,
cada um dos quais é !!a{element} de~$A$.
\end{ex}

\end{document}
